This chapter provides the context for this project and outlines the objectives to be achieved.

\section{Context}

The rapid growth of the Internet of Things (IoT) in healthcare has led to a proliferation of connected medical devices in smart hospital environments --- patient monitors, surgical equipment, diagnostic imaging systems, and environmental sensors \cite{swain2024mdafto}. These devices generate computational tasks that must be processed in real time. For instance, a surgery monitoring system cannot tolerate the 100+ millisecond latency of a distant cloud data centre; it requires single-digit millisecond response times.

Fog computing addresses this need by placing small edge servers --- called Fog Nodes (FNs) --- close to the data source, typically within the same building or campus \cite{chiti2018matching}. Compared to cloud computing, fog offers lower latency, reduced bandwidth consumption, higher reliability against internet outages, and stronger data privacy since sensitive patient information stays within the hospital network.

The core challenge addressed in this work is \textit{task offloading}: deciding which FN should process which task. Each IoT task has a strict deadline, each FN has limited capacity measured in Virtual Resource Units (VRUs), and network bandwidth is shared. Critically, not all tasks are equally important --- a major surgery matters far more than a routine checkup. If a task cannot be assigned to any FN, it results in an \textit{outage}. In healthcare, outages for critical operations are unacceptable.

This problem is fundamentally a multi-criteria matching problem: tasks have preferences over FNs (based on delay and energy), FNs have preferences over tasks (based on urgency), and both sides have constraints (deadlines, quotas). Standard load balancing optimises for a single metric and is insufficient. The Deferred Acceptance (DA) algorithm from matching theory \cite{gale1962college} provides a principled framework, and Multi-Stage Deferred Acceptance (MSDA) extends it with minimum-quota enforcement for fairness \cite{fragiadakis2016strategyproof}.

\section{Objectives}

The primary base paper, M-DAFTO \cite{swain2024mdafto}, proposes a Single-Type MSDA algorithm for IoT-fog task offloading. While effective, six key limitations can be identified:

\begin{itemize}
\item \textbf{Single-pool matching}: All tasks compete in one undifferentiated pool. A routine checkup competes for the same FN slots as a critical surgery, and under high load, critical tasks can be starved of resources.
\item \textbf{2-criteria urgency}: The urgency function uses only delay difference and preference count. Energy consumption --- a key concern in battery-powered IoT devices --- is completely ignored.
\item \textbf{No severity awareness}: There is no concept of task severity (Major vs Minor). The algorithm cannot prioritise critical healthcare operations over routine ones.
\item \textbf{Trivially consistent 2$\times$2 AHP}: With only 2 criteria, the Analytic Hierarchy Process (AHP) matrix is 2$\times$2 and always passes the consistency check (CR = 0 by definition).
\item \textbf{Delay-only normalisation}: Only delay is normalised via Min-Max scaling. Energy is neither normalised nor combined into the preference ordering.
\item \textbf{Single preference list per FN}: Each FN has one preference ranking over all tasks. There are no separate rankings for different task types.
\end{itemize}

The main objective is to adapt and implement a severity-aware 2-Type MSDA algorithm from \cite{oomori2026strategyproof} that overcomes all six limitations, while maintaining the theoretical guarantees of the original MSDA framework. Specifically, the project aims to:

\begin{enumerate}
\item Apply Major/Minor task classification with separate quota constraints based on the 2-Type framework of \cite{oomori2026strategyproof}.
\item Extend the urgency function from 2 to 4 criteria (adding energy and severity).
\item Implement a 4$\times$4 AHP matrix with genuine consistency validation.
\item Apply severity-weighted normalisation to both delay and energy.
\item Generate separate fog-side preference lists for Major and Minor tasks, as required by the 2-Type MSDA architecture \cite{oomori2026strategyproof}.
\item Evaluate the proposed method against two baselines across multiple load scales.
\end{enumerate}